% 第13章第4节：计数和求和（可行配置计数、求和、无保留变量时的汇总、保留变量的累计表、计算正确性、求和与取最小的顺序、四环算例与结果检验）。

\section{计数和求和}
\label{sec:graph-dp-semiring}
\label{sec:graph-dp-counting-sums}

最小成本只关心最好的一个配置。另一类问题要对全部配置汇总：可行配置有多少个，或者按局部权重之积给每个配置一个权重，全部权重加起来是多少。本节先用约束指示函数计算可行配置数，再把指示值推广为非负有限权重。两种计算都把“局部项相加、对变量取最小”换成“局部值相乘、对变量求和”，结构、归属和处理次序都不变；变化在于数值的组合与汇总方式，以及输出的含义。

\subsection{可行配置计数}
\label{sec:graph-dp-feasible-count}

把每个局部函数换成它的可行性指示
\[
\iota_a(\bm x_{S_a})=\begin{cases}1,&\phi_a(\bm x_{S_a})<+\infty,\\0,&\phi_a(\bm x_{S_a})=+\infty.\end{cases}
\]
一个配置可行，当且仅当这些指示全为$1$，即它们的乘积为$1$。于是可行配置数为
\begin{equation}
N=\sum_{\bm x\in\mathcal X}\ \prod_{a\in A}\iota_a(\bm x_{S_a}).
\label{eq:graph-dp-feasible-count}
\end{equation}
输入与最小成本问题相同（只用到哪些组合被禁止），输出是一个非负整数。

在袋树上计算时，节点$u$先把归入它的指示值与子节点交来的计数相乘，
\begin{equation}
h_u(\bm x_{B_u})=\prod_{a\in A_u}\iota_a(\bm x_{S_a})\cdot\prod_{c\in\operatorname{ch}(u)}n_{c\to u}(\bm x_{S_c}),
\label{eq:graph-dp-count-merge}
\end{equation}
再对接口外的状态求和，得到发往父节点的计数消息
\begin{equation}
n_{u\to p(u)}(\bm x_{S_u})=\sum_{\bm x_{I_u}}h_u(\bm x_{S_u},\bm x_{I_u}).
\label{eq:graph-dp-count-message}
\end{equation}
它的含义是：接口取这组值时，$u$下方的变量有多少种取值能使下方全部约束满足。

两种空运算的约定与最小成本不同：空乘积取$1$，求和从$0$开始。没有局部项归入、也没有子节点的袋，$h_u\equiv1$，但求和仍要累计内部变量的全部状态。孤立变量$x_v$单独成袋$\{v\}$时，消息是$\sum_{x_v}1=|\mathcal X_v|$，不能直接把它的计数设为$1$：它不受约束，每个状态都与其余变量的每个可行取值组合成不同的可行配置。

\subsection{求和}
\label{sec:graph-dp-weighted-sum}

把指示值推广为非负有限的局部权重$w_a(\bm x_{S_a})\geq0$，配置的权重定义为各局部权重之积
\[
W(\bm x)=\prod_{a\in A}w_a(\bm x_{S_a}),
\]
要求的是全部配置的权重总和
\begin{equation}
Z=\sum_{\bm x\in\mathcal X}\ \prod_{a\in A}w_a(\bm x_{S_a}).
\label{eq:graph-dp-weighted-sum}
\end{equation}
可行配置计数是取$w_a=\iota_a$的特例。

袋树上的计算与计数相同，只是把指示值换成权重。把本袋权重与子消息的乘积记为
\[
h_u(\bm x_{B_u})=\prod_{a\in A_u}w_a(\bm x_{S_a})\cdot\prod_{c\in\operatorname{ch}(u)}z_{c\to u}(\bm x_{S_c}),
\]
再对接口外状态求和，得到消息
\begin{equation}
z_{u\to p(u)}(\bm x_{S_u})=\sum_{\bm x_{I_u}}h_u(\bm x_{B_u});
\label{eq:graph-dp-sum-message}
\end{equation}
若没有接口外变量，就直接把这个乘积作为消息。

消息表照样以接口为索引，但不再需要回溯表：计数和加权总和是对全部配置的汇总，没有“达到它的配置”。

链上这一递推就是矩阵乘法。把$z_{k-1\to k}$看作行向量，把成对权重看作矩阵，矩阵的元素是$w_{k-1,k}(x_{k-1},x_k)$。于是求和
\[
\sum_{x_{k-1}}z_{k-1\to k}(x_{k-1})\,w_{k-1,k}(x_{k-1},x_k)
\]
正是行向量右乘矩阵。一般袋上对多个变量相乘再求和，则是沿指定轴汇总的张量运算。这些运算可参考\BookChapterRef{chap:linear-operations}{卷二“向量、矩阵与张量的运算”章}。

\subsection{无保留变量时的计数与求和汇总}
\label{sec:graph-dp-total-aggregation}

不保留变量时，根$r$合并得到$h_r$，再对根袋的全部状态求和：
\begin{equation}
Z=\sum_{\bm x_{B_r}}h_r(\bm x_{B_r}),
\label{eq:graph-dp-root-sum}
\end{equation}
输出一个总计数或加权总和。根袋没有变量时（例如只用来连接各分量的空袋），直接取它的合并值。

原始图不连通时，各分量的根经空接口交来标量，即各分量的总计数或总和；空接口意味着结果是单表项标量，但发出它的节点仍须对自己袋中的状态求和。这些标量作为因子乘进根的合并值，所以非连通分量的结果相乘。这与有限集合的乘法原理一致：各分量的取值互相独立地组合。

\paragraph{三元作用域例子}
表~\ref{tab:graph-dp-weights}给出三元作用域例子的权重表，权重恰在禁止的组合处为$0$。在分解$\{1,2,3\}-\{2,3,4\}$上，表~\ref{tab:graph-dp-three-aggregations}对接口的每种取值并列三种任务的左袋消息和右袋汇总值，最后一行是根处的结果：最小成本为$0$，可行配置有$10$个，权重总和为$19$。逐一枚举$16$个配置可以核对这三个数，其中最小成本和可行配置数已在第~\ref{sec:graph-dp-problem}小节的枚举中得到。

\begin{table}[!htb]
  \centering
  \caption{三元作用域例子的非负局部权重。第一列的三位二进制串对$w_A$表示$(x_1,x_2,x_3)$，对$w_B$表示$(x_2,x_3,x_4)$；权重为$0$的位置恰是表~\ref{tab:graph-dp-ternary-tables}中成本为$+\infty$的位置。}
  \label{tab:graph-dp-weights}
  \begin{tabular}{ccccccccc}
    \toprule
    取值 & $000$ & $001$ & $010$ & $011$ & $100$ & $101$ & $110$ & $111$\\
    \midrule
    $w_A$ & $1$ & $2$ & $1$ & $0$ & $2$ & $1$ & $1$ & $1$\\
    $w_B$ & $1$ & $1$ & $2$ & $0$ & $0$ & $2$ & $1$ & $2$\\
    \bottomrule
  \end{tabular}
\end{table}

\begin{table}[!htb]
  \centering
  \caption{三种任务在同一接口上的消息。最小成本使用$\min$与$+$，可行计数与加权总和使用$\sum$与$\times$。“左”列是从$\{1,2,3\}$发出的消息，对$x_1$汇总；“右”列是根袋$\{2,3,4\}$对$x_4$的汇总；根处对每个接口取值组合两列（最小成本相加，其余相乘），再对接口汇总（最小成本取最小，其余求和）。}
  \label{tab:graph-dp-three-aggregations}
  \begin{tabular}{ccccccc}
    \toprule
    & \multicolumn{2}{c}{最小成本} & \multicolumn{2}{c}{可行计数} & \multicolumn{2}{c}{加权总和}\\
    \cmidrule(lr){2-3}\cmidrule(lr){4-5}\cmidrule(lr){6-7}
    $(x_2,x_3)$ & 左 & 右 & 左 & 右 & 左 & 右\\
    \midrule
    $(0,0)$ & $0$ & $1$ & $2$ & $2$ & $3$ & $2$\\
    $(0,1)$ & $0$ & $0$ & $2$ & $1$ & $3$ & $2$\\
    $(1,0)$ & $1$ & $0$ & $2$ & $1$ & $2$ & $2$\\
    $(1,1)$ & $2$ & $1$ & $1$ & $2$ & $1$ & $3$\\
    \midrule
    根处 & \multicolumn{2}{c}{$\min(1,0,1,3)=0$} & \multicolumn{2}{c}{$4+2+2+2=10$} & \multicolumn{2}{c}{$6+6+4+3=19$}\\
    \bottomrule
  \end{tabular}
\end{table}

以接口取$(1,1)$的一行为例：$x_1=0$的一行被禁止，左袋的最小成本消息为$\min(+\infty,2)=2$，计数消息为$0+1=1$，加权消息为$0+1=1$。三种任务的计算量与表大小相同，但输出含义不同：最小成本是一个配置的值，可以回溯出达到它的配置；计数和加权总和是对全部配置的汇总。

\subsection{保留变量的累计表}
\label{sec:graph-dp-retained-aggregation}

给定保留集$R$，使用根袋覆盖$R$的查询结构（第~\ref{sec:graph-dp-query-assignment}小节），非根节点照常计算，根处只对$B_r\setminus R$求和：
\begin{equation}
Z_R(\bm x_R)=\sum_{\bm x_{B_r\setminus R}}h_r(\bm x_{B_r}).
\label{eq:graph-dp-retained-sum}
\end{equation}
结果以$\bm x_R$为索引，逐项表示把$R$固定成这组值时的可行配置数（指示权重）或权重总和。再对所有$\bm x_R$求和，才得到第~\ref{sec:graph-dp-total-aggregation}小节的总量$Z$。

与条件最小成本表相比，累计表的每一项不是一个配置的值，而是一组配置的累计；它也不需要回溯表。不含$R$的分量交来的标量作为因子乘到表的每一项上。

\subsection{计算正确性}
\label{sec:graph-dp-sum-correctness}

\paragraph{求和的三条规则}
最小成本递推所用的三条最小值规则，在求和与相乘下有完全对应的形式。设$y,z$在有限集合中取值，$h,g$取非负有限值，$c$与$y$无关，则
\[
\sum_{y,z}h(y,z)=\sum_z\sum_y h(y,z),\qquad \sum_y c\,g(y)=c\sum_y g(y),
\]
并且对取值于互不相交变量组$y_1,\ldots,y_k$的函数，
\[
\sum_{y_1,\ldots,y_k}\prod_j g_j(y_j)=\prod_j\sum_{y_j}g_j(y_j).
\]
第一条是有限求和的重排；第二条是乘法对加法的分配律；第三条由分配律展开乘积得到，展开后的每一项恰对应各$y_j$的一组取值。

\paragraph{消息的含义}
仿照定理~\ref{thm:graph-dp-bag-message}，记号同前，求和任务的消息满足
\[
z_{u\to p(u)}(\bm x_{S_u})=\sum_{\bm x_{V_{D_u}\setminus S_u}}\ \prod_{a\in A_{D_u}}w_a(\bm x_{S_a}),
\]
并且$Z=\sum_{\bm x_{B_r}}h_r(\bm x_{B_r})$。证明逐字沿用那里的拆分：变量集合$I_u$与各$V_{D_{c_j}}\setminus B_u$两两不交，并恰好组成$V_{D_u}\setminus S_u$；固定$\bm x_{B_u}$后，各子树的乘积只依赖互不相交的变量组，由第三条规则分别求和，由归纳假设得到各子消息；归入$u$的权重由第二条规则提到求和号外；最后由第一条规则对$\bm x_{I_u}$求和。

这一论证同时说明每个配置恰被累计一次。下方变量的每一组取值，唯一地分解为内部变量的一组取值和各子树中变量的一组取值；分配律展开后，每组取值恰好贡献一个乘积项$\prod_{a\in A_{D_u}}w_a$。取指示权重时，这个乘积是$1$或$0$，消息就是可行取值的个数，根处给出可行配置数。

逐项固定保留变量，同样的论证给出累计表~\eqref{eq:graph-dp-retained-sum}的正确性，全表求和得到总量。固定输入、保留集和输出要求后，仅改变待求和变量的消去次序不改变结果，因为有限求和可以任意重排；定理~\ref{thm:graph-dp-elimination-equals-messages}的对应关系也照样成立。

\paragraph{共同的代数结构}
三种任务共用同一套递推，是因为它们的运算满足同样的代数规则。每种任务都用到两种运算，分工不同：一种把同一配置的各局部项组合成这个配置的值（相加或相乘），另一种把不同取值的结果汇总成一个数（取最小或全部加起来）。消息之所以能在局部提前汇总，靠的是“先组合再汇总”与“先汇总再组合”结果相同，例如$a+\min(b,c)=\min(a+b,a+c)$、$a(b+c)=ab+ac$：与被汇总变量无关的项$a$可以留到汇总之后再组合。

把这些规则抽象出来。设集合$\mathcal S$上有两种运算：$\oplus$用于跨不同取值的汇总，$\otimes$用于在一个配置内组合各局部项。若两种运算都满足结合律和交换律，各有单位元$\bar0$与$\bar1$，满足$\bar0\otimes a=\bar0$，并且$\otimes$对$\oplus$满足分配律
\[
a\otimes(b\oplus c)=(a\otimes b)\oplus(a\otimes c),
\]
就称$(\mathcal S,\allowbreak\oplus,\allowbreak\otimes,\allowbreak\bar0,\allowbreak\bar1)$为交换\term{半环}（semiring）。

最小成本对应$(\mathbb R\cup\{+\infty\},\allowbreak\min,\allowbreak+,\allowbreak+\infty,\allowbreak0)$，此时分配律就是$a+\min(b,c)=\min(a+b,a+c)$；计数对应自然数上的$(+,\times,0,1)$；加权求和对应非负实数上的$(+,\times,0,1)$；只判断是否存在可行配置时，还可以用布尔值上的“或”与“且”。

把变量消元和袋消息中的$\min$换成$\oplus$、$+$换成$\otimes$，前面的正确性证明逐字成立，因为证明只用到了结合律、交换律和分配律。半环不要求减法或除法。

\subsection{求和与取最小的顺序}
\label{sec:graph-dp-mixed-order}

上面的换序结论只针对同一种汇总：纯最小化中可以任意交换取最小的次序，纯求和中可以任意交换求和的次序。把本节的求和与第~\ref{sec:graph-dp-min-cost}节的取最小放在一起，情况就不同了。

取一张两行两列的成本表$h(y,z)$，两行分别为$(0,1)$与$(1,0)$，即$h(0,0)=h(1,1)=0$、$h(0,1)=h(1,0)=1$。每行先对$z$求和得$1$与$1$，再对两行取最小，得到
\[
\min_y\sum_z h(y,z)=\min(1,1)=1 .
\]
反过来先对每一列在两行之间取最小，两列都得到$0$，再求和，得到
\[
\sum_z\min_y h(y,z)=0+0=0 .
\]
前者要求所有$z$共用同一个$y$，后者允许$y$随$z$改变，所以两种顺序回答的是不同的问题。

\paragraph{混合汇总}
不同变量混合使用求和与取最小（或取最大）时，须按所求结果确定运算顺序，不能直接沿用纯求和或纯最小化的换序结论。例如对表$h(x,y)$，$h(0,0)=h(1,1)=4$、$h(0,1)=h(1,0)=0$，有
\[
\max_x\sum_y h(x,y)=4,\qquad \sum_y\max_x h(x,y)=8 .
\]
这类混合问题只能按限定的次序消元：被外层汇总的变量必须最后处理，消元次序因而受到约束，所需的中间表可能比单一汇总时大。把非负权重归一化为概率还要求总和为正且有限；这些概率解释、边缘化查询以及带环图上的近似消息方法，都留到后续的概率推断内容。

\subsection{四环算例与结果检验}
\label{sec:graph-dp-four-cycle-sums}

\paragraph{约束指示表与权重表}
四环算例的约束指示表由表~\ref{tab:graph-dp-four-cycle-costs}读出：只有$\iota_{3,4}(1,1)=0$与$\iota_{1,4}(0,1)=0$，其余指示值都是$1$。另取一组非负权重，在被禁止的组合处为$0$，其余取$1$或$2$，见表~\ref{tab:graph-dp-four-cycle-weights}；这些权重是教学示例，与成本表的具体数值无关。

\begin{table}[!htb]
  \centering
  \caption{四环算例的非负局部权重。行列约定与表~\ref{tab:graph-dp-four-cycle-costs}相同；权重为$0$的位置恰是成本为$+\infty$的位置。}
  \label{tab:graph-dp-four-cycle-weights}
  \begin{tabular}{cccc}
    \toprule
    $w_{1,2}(x_1,x_2)$ & $w_{2,3}(x_2,x_3)$ & $w_{3,4}(x_3,x_4)$ & $w_{1,4}(x_1,x_4)$\\
    \midrule
    $\begin{bmatrix}1&2\\2&1\end{bmatrix}$ &
    $\begin{bmatrix}2&1\\1&1\end{bmatrix}$ &
    $\begin{bmatrix}1&2\\1&0\end{bmatrix}$ &
    $\begin{bmatrix}1&0\\1&2\end{bmatrix}$\\
    \bottomrule
  \end{tabular}
\end{table}

\paragraph{总量}
在分解$\{1,2,4\}-\{2,3,4\}$上以$\{2,3,4\}$为根。叶节点对$x_1$求和：计数消息为
\[
n(x_2,x_4)=\sum_{x_1}\iota_{1,2}\,\iota_{1,4}=(2,\ 1,\ 2,\ 1),
\]
加权消息为
\[
\begin{aligned}
z(x_2,x_4)&=\sum_{x_1}w_{1,2}\,w_{1,4}\\
&=(1\cdot1+2\cdot1,\ 1\cdot0+2\cdot2,\ 2\cdot1+1\cdot1,\ 2\cdot0+1\cdot2)\\
&=(3,4,3,2),
\end{aligned}
\]
取值按$(x_2,x_4)$的字典序排列。

根处把消息与$\phi_{2,3},\phi_{3,4}$对应的指示值或权重相乘，按$(x_2,x_3,x_4)$的字典序，计数的合并值为
\[
(2,1,2,0,2,1,2,0),
\]
求和得$N=10$。加权的合并值为
\[
(6,16,3,0,3,4,3,0),
\]
求和得$Z=35$。

用表~\ref{tab:graph-dp-four-cycle-enumeration}核对：可行配置恰有$10$个。逐个计算可行配置的权重，按字典序依次为$2,1,2,2,4,16,2,1,4,1$，总和为$35$。例如最优配置$(1,0,0,1)$的权重为
\[
w_{1,2}(1,0)\,w_{2,3}(0,0)\,w_{3,4}(0,1)\,w_{1,4}(1,1)=2\cdot2\cdot2\cdot2=16.
\]

\paragraph{保留变量的累计表}
在$R=\{1,3\}$的查询结构上，两个子节点分别对$x_2$与$x_4$求和，根袋$\{1,3\}$把两条消息相乘，不再求和。按$(x_1,x_3)$的字典序，计数消息为$(2,2,2,2)$与$(1,1,2,1)$，累计表为
\[
N_R=(2,\ 2,\ 4,\ 2);
\]
加权消息为$(4,3,5,3)$与$(1,1,5,1)$，累计表为
\[
Z_R=(4,\ 3,\ 25,\ 3).
\]
全表求和分别为
\[
2+2+4+2=10=N,\qquad 4+3+25+3=35=Z.
\]

下面逐项核对。例如$x_1=1,x_3=0$时，可行配置为
\[
 (1,0,0,0),\qquad (1,0,0,1),\qquad (1,1,0,0),\qquad (1,1,0,1),
\]
共$4$个，权重之和为$4+16+1+4=25$。

\paragraph{结果的解释与边界}
计数为零表示没有可行配置：第~\ref{sec:graph-dp-four-cycle-min}小节的不可行约束变体中，叶节点的计数消息为$(0,1,1,0)$，根处要求$x_2=x_3=x_4$时这条消息为$0$，于是$N=0$，与最小成本为$+\infty$一致。

一般的加权和为零却不直接表示不可行：权重可以在可行的组合上取$0$，例如把$w_{1,2}$全部改为$0$，则$Z=0$，而可行配置仍有$10$个。

第~\ref{sec:graph-dp-four-cycle-min}小节加入孤立变量$x_5\in\{0,1,2\}$和边$\{6,7\}$后，两个新分量经空接口交来标量。袋$\{5\}$没有局部项，$h\equiv1$，计数消息为$3$。若$\phi_{6,7}$的四项都有限，袋$\{6,7\}$的计数消息为$4$；若再给它权重表
\[
w_{6,7}=\begin{bmatrix}1&1\\2&1\end{bmatrix},
\]
加权消息为$5$。根把这些标量乘进合并值，得
\[
N=10\cdot3\cdot4=120,\qquad Z=35\cdot3\cdot5=525.
\]若错把孤立变量的计数设为$1$，结果会少乘一个因子$3$。
